
\section{The cofiber of $\eta_G$ is Eisenstein series}
\label{section:cofib_Eis}

\subsubsection{} The purpose of this section is to prove Lemma \ref{lemma:cofib_eta_Eis}. We will use $\Whit(\Bun_N)$ and not $\Whit\left(\Bunc_N \right)$.

\subsection{A compatibility}

\subsubsection{} The purpose of this subsection is to prove Proposition \ref{proposition:poinc_identification}. Let us define the map (\ref{eqn:Poinc*!2}) first. Consider the factorization
    \[
    \Bun_N \overset{r}{\longrightarrow} \Bun_N/Z_G \overset{\pi}{\longrightarrow} \Bun_G.
    \]
The action of $Z_G$ on $\Bun_N$ is trivial, so we have
    \[r_! \, [z_G] \simeq r_*.\]
Combining this with the forgetting supports map for $\pi$ gives the desired morphism
    \[\Poinc_! \, [z_G] \to \Poinc_*.\]

\subsubsection{Notation} We write
    \[
    r_{G,*} = j_* \cdot r_* \cdot k_{\Bun_G} \simeq j_* \cdot r_* \cdot \omega_{\Bun_G} [-2 \delta_G] \in \DMod\left(\Bunc_G \right).
    \]
There is a canonical map
    \[r_{G,!} \, [z_G] \to r_{G,*}.\]

\subsubsection{Proof of Proposition \ref{proposition:poinc_identification}} Recall that (\ref{eqn:Poinc*!1}) is defined by applying the forgetting supports morphisms $\pi_{2,!} \to \pi_{2, *}$ to the composite
\begin{equation}
    \label{eqn:Poinc*!intermediate}
    \left( \WW \stoximes r_{G,!} \right)^\trans [z_G]
    \to \left( \WW \shoximes r_{G,!} \right)^\trans [2 \delta_G + z_G]
    \to \left( \WW \shoximes r_{G,*} \right)^\trans [2 \delta_G].
\end{equation}
The source of this morphism is $!$-extended from the open substack
    \[
    j : \Bun_N \times \BB Z_G \mon \left( \Bun_N \!\underset{\Bun_G}{\times} \Bunc_G \right)^\trans
    \]
and the target is $*$-extended from the same locus. Therefore, (\ref{eqn:Poinc*!1}) identifies with the forgetting supports morphism for
    \[\Bun_N \times \BB Z_G \to \Bun_G\]
applied to the restriction of (\ref{eqn:Poinc*!intermediate}) along $j$. This restriction identifies with the canonical map
    \[r_! \cdot \WW \, [z_G] \to r_* \cdot \WW.\]

\subsection{The cofiber is Eisenstein series}

\subsubsection{} Let us form the diagram below.
\[
\begin{tikzcd}[
  column sep={6em,between origins},
  row sep={4em,between origins}
]
  \Bun_N \arrow[r, "r"] \arrow[d] & \Bun_N / Z_G \arrow[r, "\rho"] \arrow[d] & \Bun_N / T \arrow[r] \arrow[d, "j"] & \Bun_B \arrow[d]  \\
  \Bunc_N \arrow[r] & \Bunc_N / Z_G \arrow[r] & \Bunc_N / T \arrow[r, "i"] & \Bunc_B
\end{tikzcd}
\]
Observe that $i$ is a closed immersion because it is a base change of $\BB T \mon \Bun_T$. So the morphism
    \[\Poinc_! [z_G] \to \Poinc_*\]
is induced by the proper pushforward of
    \[
    j_! \cdot \rho_! \cdot r_!
    \overset{\alpha}{\longrightarrow} j_* \cdot \rho_! \cdot r_!
    \overset{\beta}{\longrightarrow} j_* \cdot \rho_* \cdot r_!
    \simeq j_* \cdot \rho_* \cdot r_* \, [-z_G].
    \]
along
    \[\pi : \Bunc_N/T \to \Bun_G.\]
Therefore, Proposition \ref{proposition:poinc_principal_series} is a consequence of the following two assertions:
\begin{itemize}
    \item The cofiber of $\pi_* \cdot \beta$ belongs to the image of $\Eis_{B,*}$. This will follow from Proposition \ref{proposition:Laumonesque}.
    \item The cofiber of $\pi_* \cdot \alpha$ lies in the full subcategory generated by all $\Eis_{P,*}$. This will follow from Proposition \ref{proposition:unipotent_gerbe}.
\end{itemize}

\begin{proposition}
\label{proposition:Laumonesque}
The object
  \[
  \cofib\left( \rho_! \cdot r_! \cdot \WW \to \rho_* \cdot r_! \cdot \WW \right)
  \]
lies in the image of
  \[q^! : \DMod(\Bun_T) \to \DMod(\Bun_B).\]
\end{proposition}

\subsubsection{} 
\label{subsubsection:laumon_calculation}
The mechanism responsible for this result is the following computation of G. Laumon \cite{laumon_fourier}. Consider the diagram:
\[
\begin{tikzcd}[
  column sep={5em,between origins},
  row sep={3em,between origins}
]
    & \Ga \arrow[d, "\varrho"] & \\
  \pt \arrow[r, "j"] & \Ga/\Gm & \arrow[l, "i" above] \BB\Gm
\end{tikzcd}
\]
By the contraction principle,
    \[i^* \cdot \varrho_* \cdot \AS \simeq \Gamma(\Ga/\Gm, \varrho_* \cdot \AS) \simeq 0.\]
Then one obtains
    \[\rho_* \cdot \AS \simeq j_! \cdot k[1]\]
by using base change to compute $j^! \cdot \varrho_* \cdot \AS$. Therefore, we have fiber sequences
    \[
    k\,[1] \to \varrho_! \cdot \AS \to i_* \cdot k_{\BB\Gm}
    \qquad
    i_* \cdot k_{\BB \Gm} \to \varrho_* \cdot \AS \to k \, [1]
    \]
One computes that the forgetting supports map for $\AS$ identifies with
    \[\varrho_! \cdot \AS \to i_* \cdot k_{\BB\Gm} \to \varrho_* \cdot \AS,\]
so its cofiber is an extension of $k[1]$ by itself.

\subsubsection{Proof} Recall that the Whittaker character is defined via a composite
    \[
    \Bun_N \overset{\ev}{\longrightarrow} N^\ab \simeq T_\ad^+ \overset{\mu}{\longrightarrow} \Ga.
    \]
Since the map $\ev$ is $T$-equivariant, we may form the Cartesian diagram
\[
\begin{tikzcd}[
  column sep={8em,between origins},
  row sep={4em,between origins}
]
  \Bun_N \arrow[r, "r"] \arrow[d, "\ev"] & \Bun_N \times \BB Z_G \arrow[r, "\rho"] \arrow[d] & \Bun_N / T \arrow[d, "\ev/T"]  \\
  T_\ad^+ \arrow[r] & T_\ad^+ \times \BB Z_G \arrow[r, "\varrho \times \id"] & T_\ad^+ / T_\ad \times \BB Z_G.
\end{tikzcd}
\]
Since $\ev$ is smooth, we have
\begin{align*}
    & (\rho_! \to \rho_*) \cdot r_! \cdot \WW \\
    & \qquad \simeq (\rho_! \to \rho_*) \cdot (\ev \times \id)^! \cdot (\mu^! \cdot \AS \boxtimes r_! \cdot k) \\
    & \qquad \qquad \simeq (\ev/T)^! \cdot (\varrho_! \to \varrho_*) \cdot (\mu^! \cdot \AS \boxtimes r_! \cdot k )
\end{align*}
by base change. Therefore, it suffices to prove that the object
  \[\cofib\left(\varrho_! \cdot \mu^! \cdot \AS \to \varrho_* \cdot \mu^! \cdot \AS \right)\]
lies in the full subcategory of $\DMod(T_\ad^+/T_\ad)$ generated by the constant sheaf. But $\varrho$ identifies with
  \[\Ga \times \cdots \times \Ga \to (\Ga / \Gm) \times \cdots \times (\Ga / \Gm)\]
and
  \[\mu^! \cdot \AS \simeq \AS \boxtimes \cdots \boxtimes \AS \]
because $\AS$ is a character sheaf. Therefore, the aforementioned calculation (\S\ref{subsubsection:laumon_calculation}) induces a filtration on the cofiber in question by the power poset of $\Delta$, each of whose graded pieces is
    \[k[1] \boxtimes \cdots \boxtimes k[1].\]

\begin{proposition}
\label{proposition:unipotent_gerbe}
The direct image along the proper morphism $\Bunc_N/T \to \Bun_G$ of any object supported on the complement of
    \[\Bun_N/T \mon \Bunc_N/T\]
is Eisenstein series.
\end{proposition}

\subsubsection{} In what follows, we will prove this proposition in several steps under the assumption that $X$ does not have genus zero. In the genus zero case, the entire category $\DMod(\Bun_G)$ is Eisenstein series if $G$ is not a torus. The proposition is vacuous when $G$ is a torus.

\subsubsection{First step} To analyse this complement, recall that $\Bunc_B$ has a defect stratification indexed by dominant cocharacters
  \[
  \left(X^\lambda \times \Bun_T \right) \underset{\Bun_T}{\times} \Bun_B \simeq \Bunc_B^\lambda \mon \Bunc_B.
  \]
Here we are using the map
    \[
    X^\lambda \times \Bun_T \to \Bun_T 
    \qquad (D, \FF_T) \mapsto \FF_T(D).
    \]
We will use the notation $\Bunc_N^\lambda / T$ for the induced stratification of $\Bunc_N / T$. It suffices to prove that for every nonzero $\lambda$, the image of $*$-pushforward along the map
    \[\Bunc_N^\lambda/T \to \Bun_G\]
is Eisenstein series.

\subsubsection{Second step} Let $P$ be an arbitrary standard parabolic. Then we have the following Cartesian diagram:
\[
\begin{tikzcd}[
  column sep={5.5em,between origins},
  row sep={4em,between origins}
]
  \Bunc_B^\lambda \arrow[r] \arrow[d] & \Bun_B \arrow[r] \arrow[d] & \Bun_P \arrow[d] \\
  \Bunc_{B_M}^\lambda \arrow[r] \arrow[d] & \Bun_{B_M} \arrow[r] \arrow[d] & \Bun_M \\
  X^\lambda \times \Bun_T \arrow[r] & \Bun_T &
\end{tikzcd}
\]
This gives the Cartesian left face of the following diagram:
\[
\begin{tikzcd}[
  column sep={6em,between origins},
  row sep={2em,between origins}
]
   & \Bunc_N^\lambda / T \arrow[r] \arrow[dl, " 'q" above left] \arrow[dd] & \Bunc_N / T \arrow[ddd, bend left=10]  \\
  \Bunc_{N_M}^\lambda / T \arrow[dd] & \\
  & \Bun_P \arrow[dl] \arrow[dr] & \\
  \Bun_M & & \Bun_G
\end{tikzcd}
\]
By base change, it suffices to show that for every nonzero $\lambda$, there exists a proper parabolic for which the image of the functor
  \[
  {'q}^! : \DMod\left( \Bunc_{N_M}^\lambda / T \right) \to \DMod\left( \Bunc_N^\lambda / T \right)
  \]
generates the target.

\subsubsection{Third step} For each $\lambda$, let $P_\lambda$ be the parabolic corresponding to the set of positive roots $\alpha$ for which $\lambda(\alpha) \neq 0$. We claim that if the genus of $X$ is nonzero, then after taking $P = P_\lambda$ in the above notation, $'q$ is tower of unipotent gerbes. This will prove the proposition because $P_\lambda$ is a proper parabolic whenever $\lambda \neq 0$.

\subsubsection{Facts} Let us recall some basic facts before proving the claim.
\begin{itemize}
    \item[(i)] Let $H \to G$ be a map of groups. Then the fiber of $\Bun_H \to \Bun_G$ over a point $\FF_G$ classifies sections of \[\FF_G \overset{G}{\times} G/H.\]
    \item[(ii)] Let $\Gm$ act on $\BB \Ga$ through the dilation action on $\Ga$. For any line bundle $\LL$ on $X$, sections of $\LL \overset{\Gm}{\times} \BB \Ga$ are classified by $\Gamma(X, \LL[1])$.
\end{itemize}
Note that (ii) is obtained from (i) by taking $H \mon \GL_2$ to be the automorphism group of $\AA^1$.

\subsubsection{Notation} Choose a sequence of quotients
  \[B = B_0 \epi B_1 \epi \cdots \epi B_n = B/U_P\]
with the property that
  \[\ker(B_i \epi B_{i+1}) \simeq \Ga\]
is central in the unipotent radical of $B_i$. Then the action of $B_{i+1}$ on $B_{i+1}/B_i \simeq \BB \Ga$ factors through $B_{i+1} \epi T$. The weight of the $T$-action is some positive root $\alpha_i$ of $M$.

\subsubsection{Proof of the claim} It suffices to show that the map $\aph q_i$ in the Cartesian diagram below is a unipotent gerbe.
\[
\begin{tikzcd}[
  column sep={6em,between origins},
  row sep={4em,between origins}
]
  \YY_i \arrow[r, "\aph q_i"] \arrow[d] & \YY_{i+1} \arrow[r] \arrow[d] & X^\lambda \arrow[d, "\AJ_\omega"] \\
  \Bun_{B_i} \arrow[r] & \Bun_{B_{i+1}} \arrow[r] & \Bun_T
\end{tikzcd}
\]
So fix a point of $\YY_{i+1}$, and say that its image in $X^\lambda$ is $D$. Then the fiber of $\aph q_i$ over our chosen point is
\begin{equation}
\label{eqn:unipotent_gerbe_complex}
    \Gamma\left(X, \omega^{\alpha_i(\rhov)}(\alpha_i(D))[1] \right)
    \simeq \Gamma \left(X, \omega^{1 - \alpha_i(\rhov)} (-\alpha_i(D)) \right)^\vee
\end{equation}
We claim that this complex is concentrated in negative degrees. Indeed, the line bundle appearing in the RHS of this expression has degree
    \[
    (1 - \alpha_i(\rhov)) \cdot \deg(\omega_X) - \alpha_i(D).
    \]
This integer is negative because $\alpha_i(\rhov)$ is a positive integer, $\alpha_i(D)$ is a positive divisor of degree $\alpha_i(\lambda) > 0$, and $\deg(\omega_X) \ge 0$ by our assumption on the genus.
